\documentclass[11pt, a4paper]{article}

\usepackage[T1]{fontenc}
\usepackage[utf8]{inputenc}
\usepackage{tgpagella}
\usepackage{microtype}
\usepackage[spanish, es-noshorthands]{babel}
\usepackage{amsmath, amssymb}
\usepackage[margin=2.5cm]{geometry}
\usepackage{fancyhdr}
\usepackage{tcolorbox}

\pagestyle{fancy}
\fancyhf{}
\setlength{\headheight}{16pt}
\lhead{\textbf{IMT-342} | Robótica}
\chead{Puente Curricular (Provisional)}
\rhead{Semana 6}

\begin{document}

\begin{center}
    \Large\textbf{De la Pose al Movimiento: Limitaciones de la Cinemática Directa}\\[0.2cm]
    \normalsize \textit{Nota: Material provisional pendiente de confirmación de syllabus[cite: 8].}
\end{center}

\section*{¿Por qué la Cinemática Directa no es suficiente?[cite: 8]}
La cinemática directa que acabamos de validar responde a la pregunta de \textbf{posición}:
\begin{equation*}
    q \longmapsto x(q)
\end{equation*}
donde $x$ representa la pose del efector final. Sin embargo, si el robot se está moviendo con velocidades articulares $\dot{q}$, ¿qué tan rápido y en qué dirección se mueve el efector final?

Por regla de la cadena multivariable para las coordenadas de posición $p = p(q)$:
\begin{equation*}
    \dot{p} = \frac{\partial p}{\partial q}\dot{q}
\end{equation*}

Definimos provisionalmente al Jacobiano $J_p(q) = \frac{\partial p}{\partial q}$, de modo que:
\begin{equation*}
    \dot{p} = J_p(q)\dot{q}
\end{equation*}

\section*{Demostración Preliminar 2R Planar[cite: 8]}
Para un brazo 2R:
\begin{align*}
    p_x &= a_1\cos q_1 + a_2\cos(q_1+q_2) \\
    p_y &= a_1\sin q_1 + a_2\sin(q_1+q_2)
\end{align*}

Derivando con respecto al tiempo:
\begin{align*}
    \dot{p}_x &= [-a_1\sin q_1 - a_2\sin(q_1+q_2)]\dot{q}_1 - a_2\sin(q_1+q_2)\dot{q}_2 \\
    \dot{p}_y &= [a_1\cos q_1 + a_2\cos(q_1+q_2)]\dot{q}_1 + a_2\cos(q_1+q_2)\dot{q}_2
\end{align*}

Escrito en forma matricial:
\begin{equation*}
    \begin{bmatrix} \dot{p}_x \\ \dot{p}_y \end{bmatrix} = 
    \underbrace{\begin{bmatrix} -a_1s_1 - a_2s_{12} & -a_2s_{12} \\ a_1c_1 + a_2c_{12} & a_2c_{12} \end{bmatrix}}_{J_p(q)}
    \begin{bmatrix} \dot{q}_1 \\ \dot{q}_2 \end{bmatrix}
\end{equation*}

\begin{tcolorbox}[colback=white, colframe=black, title=Pregunta de Transición]
    \textit{La cinemática directa proporciona la pose a partir de las posiciones articulares. ¿Qué nuevo mapeo (matriz) necesitamos para predecir la velocidad del efector final a partir de las velocidades articulares?} \\
    \textbf{Respuesta esperada:} Una matriz de mapeo de velocidades o Matriz Jacobiana.
\end{tcolorbox}

\end{document}
